\intro

Современное развитие техники и технологий сопровождается ростом требований к точности расчёта тепловых процессов, протекающих в различных материалах и конструкциях. Анализ теплопередачи~\cite{krainov2016} является важной задачей при проектировании изделий в машиностроении, электронике, энергетике и других отраслях, где температурные режимы оказывают существенное влияние на работоспособность и надёжность систем.

В условиях усложнения геометрии объектов и физических условий эксплуатации аналитические методы решения задач теплопроводности оказываются ограниченными в применении. В связи с этим широкое распространение получили численные методы, позволяющие исследовать нестационарные процессы теплопередачи в многомерных областях с заданными начальными и граничными условиями.

Для решения задач теплопередачи на практике широко применяются специализированные программные комплексы, среди которых одним из наиболее распространённых является «COMSOL Multiphysics»~\cite{vajdi2020review}. Данный программный пакет реализует численное моделирование на основе метода конечных элементов~\cite{розин2000метод} и предоставляет широкие возможности для анализа тепловых процессов в сложных геометрических областях.

Использование «COMSOL Multiphysics»~\cite{suarez2014study} позволяет задавать физические свойства материалов, начальные и граничные условия, а также получать наглядные результаты в виде распределений температуры и тепловых потоков. Это делает его удобным инструментом для инженерных расчётов и научных исследований.

Вместе с тем применение готовых программных комплексов имеет ряд ограничений, связанных с высокой вычислительной сложностью, избыточностью функционала для узкоспециализированных задач, а также ограниченной гибкостью при реализации собственных численных алгоритмов.

В связи с этим целью данной работы является разработка программного комплекса для численного моделирования нестационарной теплопередачи в сплошных средах на основе уравнения теплопроводности Фурье–Кирхгофа.

Для достижения поставленной цели необходимо решить следующие задачи:
\begin{enumerate}[label=\arabic*.]
    \item провести анализ физических основ теплопроводности и математической модели нестационарной теплопередачи;
    \item выполнить постановку задачи с заданием начальных и граничных условий в трёхмерной области;
    \item реализовать численный алгоритм решения уравнения теплопроводности на основе явной конечно-разностной схемы;
    \item разработать программный комплекс HeatSim3D, обеспечивающий проведение расчётов;
    \item реализовать средства визуализации и анализа температурных полей;
    \item провести вычислительные эксперименты и выполнить анализ полученных результатов.
\end{enumerate}

В первой главе рассматриваются физические основы теплопроводности, приводится вывод уравнения теплопередачи и формулируется математическая постановка задачи. Описывается конечно-разностный метод решения и анализируются его свойства.

Во второй главе представлена разработка программного комплекса, включая описание архитектуры, реализацию численного алгоритма, организацию вычислительного процесса и средства визуализации результатов.

В третьей главе приведены результаты численного моделирования, выполнена верификация разработанной модели и проведён анализ влияния параметров на распределение температуры.

